\setcounter{section}{7}
\setcounter{theorem}{0}
\setcounter{definition}{0}
\setcounter{remark}{0}
\setcounter{note}{0}

\subsubsection{Preliminaries, then the script}

The problem is unconstrained convex minimization of
\begin{equation}
\label{eq:p7-primal}
f(x)
\coloneqq
\|Ax-b\|_{2}+\gamma\|x\|_{1},
\qquad
x\in\mathbb{R}^{n}.
\end{equation}
The \(\ell_1\) term is not differentiable on the coordinate hyperplanes, and \(\|Ax-b\|_{2}\) is not differentiable on \(\{Ax=b\}\). Part~(a) therefore rewrites the residual as an equality constraint and passes to the Lagrange dual, which is a constrained problem in a multiplier \(\nu\in\mathbb{R}^{m}\) with only ball constraints. Parts~(b) and~(c) are first-order optimality for \(f\) itself. The dictionary is short; each item is used on the next page.

\begin{definition}[The three norms]
\label{def:p7-lp}
\[
\|x\|_{1}=\sum_{i=1}^{n}|x_{i}|,
\qquad
\|x\|_{2}=\Bigl(\sum_{i=1}^{n}x_{i}^{2}\Bigr)^{1/2},
\qquad
\|x\|_{\infty}=\max_{i}|x_{i}|.
\]
The only two pairings needed below are Cauchy--Schwarz and H\"older:
\begin{equation}
\label{eq:p7-holder}
\lvert\nu^{\top}y\rvert\le\|\nu\|_{2}\|y\|_{2},
\qquad
\lvert u^{\top}x\rvert
=
\Bigl\lvert\sum_{i}u_{i}x_{i}\Bigr\rvert
\le
\sum_{i}\lvert u_{i}\rvert\lvert x_{i}\rvert
\le
\|u\|_{\infty}\|x\|_{1}.
\end{equation}
If \(a_{i}\) is the \(i\)-th column of \(A\), then \((A^{\top}r)_{i}=a_{i}^{\top}r\).
\end{definition}

\begin{definition}[Lagrangian and Lagrange dual]
\label{def:p7-dual}
For \(\min_{z}F(z)\) subject to \(h(z)=0\), the Lagrangian is \(L(z,\nu)=F(z)+\nu^{\top}h(z)\). The dual function is \(g(\nu)=\inf_{z}L(z,\nu)\), and the Lagrange dual problem is \(\sup_{\nu}g(\nu)\). Weak duality is the comparison \(g(\nu)\le F(z)\) for every feasible \(z\) and every \(\nu\). Here \(z=(x,y)\), \(F=\|y\|_{2}+\gamma\|x\|_{1}\), and the exam's constraint \(Ax-b=y\) is written \(h(x,y)=Ax-b-y\).
\end{definition}

\begin{lemma}[When \(\inf_{z}(\alpha\|z\|-u^{\top}z)\) is finite]
\label{lem:p7-norm-inf}
Let \(\alpha>0\). Then
\[
\inf_{y}\bigl(\|y\|_{2}-\nu^{\top}y\bigr)
=
\begin{cases}
0, & \|\nu\|_{2}\le 1,\\
-\infty, & \|\nu\|_{2}>1,
\end{cases}
\qquad
\inf_{x}\bigl(\gamma\|x\|_{1}+u^{\top}x\bigr)
=
\begin{cases}
0, & \|u\|_{\infty}\le\gamma,\\
-\infty, & \|u\|_{\infty}>\gamma.
\end{cases}
\]
\end{lemma}

\begin{proof}
If \(\|\nu\|_{2}\le 1\), then \(\|y\|_{2}-\nu^{\top}y\ge\|y\|_{2}-\|\nu\|_{2}\|y\|_{2}\ge 0\), and the value \(0\) is attained at \(y=0\). If \(\|\nu\|_{2}>1\), the ray \(y=t\nu\) with \(t\to+\infty\) gives \(t\|\nu\|_{2}-t\|\nu\|_{2}^{2}=t\|\nu\|_{2}(1-\|\nu\|_{2})\to-\infty\).

If \(\|u\|_{\infty}\le\gamma\), then \(\gamma\|x\|_{1}+u^{\top}x\ge\gamma\|x\|_{1}-\|u\|_{\infty}\|x\|_{1}\ge 0\), attained at \(x=0\). If \(\|u\|_{\infty}>\gamma\), some coordinate satisfies \(\lvert u_{j}\rvert>\gamma\). The ray \(x=-t\,\operatorname{sign}(u_{j})\,e_{j}\) gives \(\gamma t-t\lvert u_{j}\rvert=t(\gamma-\lvert u_{j}\rvert)\to-\infty\).
\end{proof}

\subsubsection{(a) The Lagrange dual}

The original \(f(x)\) and the constrained problem in the prompt are the same object: given any \(x\), the only feasible \(y\) is \(y=Ax-b\), and the objective becomes \(f(x)\). The reason for introducing \(y\) is that the two norms then act on separate variables, so the dual function splits.

Form the Lagrangian of the written constraint \(Ax-b-y=0\):
\begin{equation}
\label{eq:p7-L}
L(x,y,\nu)
=
\|y\|_{2}+\gamma\|x\|_{1}+\nu^{\top}(Ax-b-y).
\end{equation}
The dual function is the unconstrained infimum
\[
g(\nu)
=
\inf_{x\in\mathbb{R}^{n},\,y\in\mathbb{R}^{m}}L(x,y,\nu).
\]
There is no \(x\)--\(y\) product in \(L\) except through the already-fixed vector \(\nu\), so the infimum factors:
\begin{equation}
\label{eq:p7-g-split}
g(\nu)
=
\inf_{y}\bigl(\|y\|_{2}-\nu^{\top}y\bigr)
+
\inf_{x}\bigl(\gamma\|x\|_{1}+(A^{\top}\nu)^{\top}x\bigr)
-
b^{\top}\nu.
\end{equation}
Apply \cref{lem:p7-norm-inf} to each infimum (the second with \(u=A^{\top}\nu\)). Both are \(0\) precisely when
\[
\|\nu\|_{2}\le 1
\qquad\text{and}\qquad
\|A^{\top}\nu\|_{\infty}\le\gamma;
\]
otherwise at least one infimum is \(-\infty\), hence \(g(\nu)=-\infty\). On that set one is left with the linear function \(-b^{\top}\nu\). The Lagrange dual is therefore the explicit program
\begin{equation}
\label{eq:p7-dual-prog}
\sup_{\nu\in\mathbb{R}^{m}}\bigl(-b^{\top}\nu\bigr)
\qquad\text{subject to}\qquad
\|\nu\|_{2}\le 1,
\quad
\|A^{\top}\nu\|_{\infty}\le\gamma.
\end{equation}
Equivalently, after the change of sign \(\lambda=-\nu\) (which is the Lagrangian of the flipped writing \(y-Ax+b=0\)),
\begin{equation}
\label{eq:p7-dual-flip}
\sup_{\lambda\in\mathbb{R}^{m}}\,b^{\top}\lambda
\qquad\text{subject to}\qquad
\|\lambda\|_{2}\le 1,
\quad
\|A^{\top}\lambda\|_{\infty}\le\gamma.
\end{equation}
Either display is a correct answer. Weak duality is \(g(\nu)\le f(x)\) for every \(x\) and every feasible \(\nu\) (\cref{rem:p7-weak}). Strong duality holds because the constraint is affine and a feasible pair always exists (pick any \(x\), set \(y=Ax-b\)): the optimal values of \eqref{eq:p7-primal} and \eqref{eq:p7-dual-prog} coincide (\cref{rem:p7-strong}).

\begin{remark}[Weak duality]
\label{rem:p7-weak}
Write
\[
p^{*}
\coloneqq
\inf_{x\in\mathbb{R}^{n}}f(x)
=
\inf_{\substack{x,y\\Ax-b=y}}
\bigl(\|y\|_{2}+\gamma\|x\|_{1}\bigr)
\]
for the primal value, and
\[
d^{*}
\coloneqq
\sup_{\nu\in\mathbb{R}^{m}}g(\nu)
\]
for the dual value. Weak duality is the comparison of these two numbers through a pointwise inequality that does not mention optima.

Take any \(x\in\mathbb{R}^{n}\) and set \(y=Ax-b\), so that \((x,y)\) is feasible for the written constraint. The Lagrangian \eqref{eq:p7-L} then collapses to the primal objective:
\[
L(x,y,\nu)
=
\|y\|_{2}+\gamma\|x\|_{1}+\nu^{\top}(Ax-b-y)
=
f(x).
\]
The dual function is an unconstrained infimum over a larger set, hence cannot exceed any particular value of \(L\):
\[
g(\nu)
=
\inf_{x',y'}L(x',y',\nu)
\le
L(x,y,\nu)
=
f(x).
\]
This holds for every \(x\) and every multiplier \(\nu\), with no convexity assumption and with no restriction that \(\nu\) be dual-feasible \cite[\S5.1.3]{boyd04}. Taking the infimum in \(x\) on the right and the supremum in \(\nu\) on the left yields
\[
d^{*}\le p^{*}.
\]
The difference \(p^{*}-d^{*}\ge 0\) is the duality gap.

If \(\nu\) fails the two ball constraints of \eqref{eq:p7-dual-prog}, then \(g(\nu)=-\infty\) by \cref{lem:p7-norm-inf}, and the inequality is vacuous. If \(\nu\) is dual-feasible, then \(g(\nu)=-b^{\top}\nu\), and weak duality becomes a concrete numerical bound: every such \(\nu\) lower-bounds every Lasso value,
\[
-b^{\top}\nu
\le
f(x)
\qquad
\text{for all }x\in\mathbb{R}^{n}.
\]
In particular \(-b^{\top}\nu\le p^{*}\). That is the whole point of writing a dual even if one never proves that the gap is zero: any feasible \(\nu\) certifies a lower bound, and a better \(\nu\) certifies a tighter one.

Weak duality does not say that \(g(\nu)=f(x)\) at some pair, and it does not say that a minimizer of \(f\) exists. It is only the one-sided comparison \(g\le f\).
\end{remark}

\begin{remark}[Strong duality]
\label{rem:p7-strong}
Strong duality is the statement that the gap vanishes:
\[
d^{*}=p^{*}.
\]
The best lower bound produced by the dual is then exactly the primal value, so the two programs in \eqref{eq:p7-primal} and \eqref{eq:p7-dual-prog} have the same optimal value. This is not implied by weak duality, and it is not implied by convexity of \(f\) alone.

A standard sufficient condition is Slater's constraint qualification for convex programs. In the form needed here \cite[\S5.2.3]{boyd04}: the objective \(F(x,y)=\|y\|_{2}+\gamma\|x\|_{1}\) is convex, the only constraint \(Ax-b-y=0\) is affine, and there exists a feasible point in the relative interior of \(\operatorname{dom} F\). Under those hypotheses one has \(d^{*}=p^{*}\).

Each hypothesis is immediate. The function \(F\) is a sum of norms, hence convex, and it is finite at every \((x,y)\in\mathbb{R}^{n}\times\mathbb{R}^{m}\), so \(\operatorname{ri}(\operatorname{dom} F)\) is the whole space. The constraint is affine by inspection. A feasible pair always exists: pick any \(x\) and set \(y=Ax-b\). Slater therefore applies, and the optimal values of \eqref{eq:p7-primal} and \eqref{eq:p7-dual-prog} coincide.

(The same conclusion is Fenchel--Rockafellar applied to \(\|A\,\cdot\,-b\|_{2}+\gamma\|\cdot\|_{1}\): the conjugate of \(\|\cdot\|_{2}\) is the indicator of the Euclidean unit ball, the conjugate of \(\gamma\|\cdot\|_{1}\) is the indicator of the \(\ell_{\infty}\) ball of radius \(\gamma\), and both domains are the whole space.)

Attainment is a separate claim. Strong duality equates two numbers; it does not by itself produce a minimizer \(x^{*}\) or a maximizer \(\nu^{*}\). Both happen to exist here. The primal attains because \(f\) is continuous and coercive: \(\gamma>0\) gives \(f(x)\ge\gamma\|x\|_{1}\to\infty\) as \(\|x\|\to\infty\). The dual attains because \(-b^{\top}\nu\) is continuous and the feasible set
\[
\bigl\{\nu:\|\nu\|_{2}\le 1\bigr\}
\cap
\bigl\{\nu:\|A^{\top}\nu\|_{\infty}\le\gamma\bigr\}
\]
is compact (closed subset of a Euclidean ball). Thus there exist \(x^{*}\) and \(\nu^{*}\) with
\[
f(x^{*})=g(\nu^{*})=p^{*}=d^{*}.
\]
A pair with zero gap is a saddle of \(L\): writing \(y^{*}=Ax^{*}-b\),
\[
L(x^{*},y^{*},\nu)
\le
L(x^{*},y^{*},\nu^{*})
\le
L(x,y,\nu^{*})
\]
for every \((x,y)\) and every \(\nu\). The left inequality is primal feasibility (the constraint term vanishes at \((x^{*},y^{*})\), independently of \(\nu\)); the right inequality says that \(\nu^{*}\) attains the infimum that defines \(g\). That saddle reading is the route of \cref{rem:p7-from-dual}. Part~(b) does not need it.

Two statements that are easy to swap. Weak duality is \(g(\nu)\le f(x)\) for every pair, including non-optimal ones. Strong duality is \(d^{*}=p^{*}\), i.e.\ equality of the two optimal values, not equality of \(g\) and \(f\) at a random feasible pair. Equality \(g(\nu)=f(x)\) holds if and only if \(x\) is primal-optimal and \(\nu\) is dual-optimal (and then the gap at that pair is automatically zero).
\end{remark}

\begin{examnote}[Part~(a) on a script]
\label{examnote:p7-dual}
Write \eqref{eq:p7-L} and \eqref{eq:p7-g-split}. Evaluate the two infima by Cauchy--Schwarz and H\"older as in \cref{lem:p7-norm-inf} (the two rays that produce \(-\infty\) must be exhibited). State \eqref{eq:p7-dual-prog} or \eqref{eq:p7-dual-flip}. One sentence on weak duality (\(g\le f\)) and one on Slater (affine constraint, pick \(y=Ax-b\)) is enough; \cref{rem:p7-weak,rem:p7-strong} are not the script. Naming ``the Lagrange dual'' without the constraints on \(\nu\) is not a derivation.
\end{examnote}

\subsubsection{(b) First-order conditions at \(x^{*}\)}

Let \(x^{*}\) minimize \(f\), and assume \(y^{*}\coloneqq Ax^{*}-b\neq 0\). The exam's vector
\[
r
=
\frac{y^{*}}{\|y^{*}\|_{2}}
\]
is the Euclidean unit vector in the direction of the residual. It is also the gradient of \(\|\cdot\|_{2}\) at \(y^{*}\), which is the only reason it appears.

Because \(x^{*}\) is an unconstrained minimizer, every directional derivative of \(f\) at \(x^{*}\) is nonnegative: for every \(v\in\mathbb{R}^{n}\),
\begin{equation}
\label{eq:p7-dir}
\liminf_{t\downarrow 0}\frac{f(x^{*}+tv)-f(x^{*})}{t}
\ge 0.
\end{equation}
Compute the two pieces of \(f\) separately.

\emph{The residual term.} Write \(y(t)=y^{*}+tAv\). Then
\[
\|y(t)\|_{2}^{2}
=
\|y^{*}\|_{2}^{2}+2t\,(y^{*})^{\top}Av+t^{2}\|Av\|_{2}^{2},
\]
and the first-order expansion of the square root about \(\|y^{*}\|_{2}>0\) is
\begin{equation}
\label{eq:p7-res-diff}
\|y^{*}+tAv\|_{2}
=
\|y^{*}\|_{2}+t\,r^{\top}Av+o(t).
\end{equation}
(This is the ordinary derivative of \(\|\cdot\|_{2}\) at a nonzero point: \(\nabla\|y\|_{2}=y/\|y\|_{2}\).)

\emph{The \(\ell_1\) term.} For all sufficiently small \(\lvert t\rvert\),
\begin{equation}
\label{eq:p7-l1-diff}
\|x^{*}+tv\|_{1}-\|x^{*}\|_{1}
=
t\sum_{x_{i}^{*}\neq 0}\operatorname{sign}(x_{i}^{*})\,v_{i}
+
\lvert t\rvert\sum_{x_{i}^{*}=0}\lvert v_{i}\rvert.
\end{equation}
(On a coordinate with \(x_{i}^{*}\neq 0\), \(\lvert x_{i}^{*}+tv_{i}\rvert=\lvert x_{i}^{*}\rvert+t\operatorname{sign}(x_{i}^{*})v_{i}\) for small \(t\). On a coordinate with \(x_{i}^{*}=0\), one has \(\lvert tv_{i}\rvert\).)

\emph{Coordinatewise tests.} Take \(v=e_{i}\) and \(v=-e_{i}\). Write \(u\coloneqq A^{\top}r\), so that \(r^{\top}Ae_{i}=u_{i}=a_{i}^{\top}r\).

If \(x_{i}^{*}\neq 0\), both directions are admissible and \eqref{eq:p7-l1-diff} is linear in \(t\). Condition \eqref{eq:p7-dir} in both directions forces the derivative to vanish:
\begin{equation}
\label{eq:p7-active}
u_{i}+\gamma\operatorname{sign}(x_{i}^{*})=0,
\qquad\text{i.e.}\qquad
\lvert u_{i}\rvert=\gamma.
\end{equation}

If \(x_{i}^{*}=0\), the \(\ell_1\) contribution is \(\gamma\lvert t\rvert\). The right derivative of \(f\) along \(+e_{i}\) is \(u_{i}+\gamma\), and the right derivative along \(-e_{i}\) is \(-u_{i}+\gamma\). Both are \(\ge 0\), hence
\begin{equation}
\label{eq:p7-inactive}
\lvert u_{i}\rvert\le\gamma.
\end{equation}

Combining \eqref{eq:p7-active} and \eqref{eq:p7-inactive} over every coordinate,
\begin{equation}
\label{eq:p7-inf}
\|A^{\top}r\|_{\infty}
=
\|u\|_{\infty}
\le\gamma,
\end{equation}
which is the first claim. For the second claim, pair \(u\) against \(x^{*}\). On every coordinate with \(x_{i}^{*}=0\) the product \(u_{i}x_{i}^{*}\) is zero. On every coordinate with \(x_{i}^{*}\neq 0\), equation \eqref{eq:p7-active} gives \(u_{i}x_{i}^{*}=-\gamma\lvert x_{i}^{*}\rvert\). Summing,
\begin{equation}
\label{eq:p7-comp}
r^{\top}Ax^{*}+\gamma\|x^{*}\|_{1}
=
u^{\top}x^{*}+\gamma\|x^{*}\|_{1}
=
0.
\end{equation}

\begin{remark}[The same argument in subdifferential language]
\label{rem:p7-subgrad}
Convexity of \(f\) says that \(x^{*}\) minimizes \(f\) if and only if \(0\in\partial f(x^{*})\). The chain rule and the sum rule give
\[
\partial f(x^{*})
=
A^{\top}\partial\|\cdot\|_{2}(y^{*})+\gamma\,\partial\|x^{*}\|_{1}.
\]
For \(y^{*}\neq 0\) one has \(\partial\|y^{*}\|_{2}=\{r\}\). The set \(\partial\|x\|_{1}\) is exactly the collection of vectors \(s\) with \(\|s\|_{\infty}\le 1\) and \(s_{i}=\operatorname{sign}(x_{i})\) whenever \(x_{i}\neq 0\); in particular \(s^{\top}x=\|x\|_{1}\). Thus \(A^{\top}r+\gamma s=0\) for some such \(s\), which is \eqref{eq:p7-inf}--\eqref{eq:p7-comp} again. This is 2024 Fall Problem~7(ii). The directional-derivative argument above does not require the vocabulary.
\end{remark}

\begin{remark}[From the dual of (a)]
\label{rem:p7-from-dual}
The same two displays are dual feasibility plus zero gap. Differentiability of \(\|y\|_{2}\) at \(y^{*}\neq 0\) forces the optimal multiplier of \eqref{eq:p7-L} to be \(\nu^{*}=r\). Dual feasibility of \(r\) is \(\|r\|_{2}=1\le 1\) and \(\|A^{\top}r\|_{\infty}\le\gamma\). Strong duality (\cref{rem:p7-strong}) then reads \(\|y^{*}\|_{2}+\gamma\|x^{*}\|_{1}=-b^{\top}r\). Substitute \(\|y^{*}\|_{2}=r^{\top}y^{*}=r^{\top}Ax^{*}-b^{\top}r\) to recover \eqref{eq:p7-comp}. This route uses (a); the directional argument does not.
\end{remark}

\subsubsection{(c) Small columns cannot be used}

The claim does not assume \(Ax^{*}\neq b\). The shortest argument uses only the triangle inequality, and covers both cases.

Suppose \(x_{i}^{*}\neq 0\). For \(0<t<\lvert x_{i}^{*}\rvert\) set \(\sigma=\operatorname{sign}(x_{i}^{*})\) and \(x(t)=x^{*}-t\sigma e_{i}\). Then the \(i\)-th coordinate of \(x(t)\) has absolute value \(\lvert x_{i}^{*}\rvert-t\), and no other coordinate changes, so
\[
\|x(t)\|_{1}=\|x^{*}\|_{1}-t.
\]
The residual becomes \(y^{*}-t\sigma a_{i}\), hence
\[
\|Ax(t)-b\|_{2}
\le
\|y^{*}\|_{2}+t\|a_{i}\|_{2}.
\]
Adding \(\gamma\|x(t)\|_{1}\),
\[
f\bigl(x(t)\bigr)
\le
f(x^{*})+t\bigl(\|a_{i}\|_{2}-\gamma\bigr).
\]
The hypothesis \(\|a_{i}\|_{2}<\gamma\) makes the right-hand side strictly smaller than \(f(x^{*})\) for every small \(t>0\), contradicting optimality of \(x^{*}\). Therefore \(x_{i}^{*}=0\).

If one has already written (b) and is in the case \(y^{*}\neq 0\), a shorter sentence is available: \(x_{i}^{*}\neq 0\) would force \(\lvert a_{i}^{\top}r\rvert=\gamma\) by \eqref{eq:p7-active}, while Cauchy--Schwarz gives \(\lvert a_{i}^{\top}r\rvert\le\|a_{i}\|_{2}\|r\|_{2}=\|a_{i}\|_{2}<\gamma\).

\begin{examnote}[Parts~(b) and~(c) on a script]
\label{examnote:p7-bc}
For (b): expand \(f(x^{*}+te_{i})\) using \eqref{eq:p7-res-diff} and \eqref{eq:p7-l1-diff}, treat \(x_{i}^{*}\neq 0\) and \(x_{i}^{*}=0\) separately, and sum \(u_{i}x_{i}^{*}\). For (c): the triangle argument is one paragraph and does not need \(r\). Do not quote ``soft thresholding'' without writing \(\lvert u_{i}\rvert\le\gamma\).
\end{examnote}

\begin{summary}[Problem~7]
\label{sum:p7}
The Lagrange dual of the residual form is \eqref{eq:p7-dual-prog} (or \eqref{eq:p7-dual-flip}). At a minimizer with nonzero residual, \(A^{\top}r\) lies in the \(\ell_{\infty}\) ball of radius \(\gamma\) and is complementary to \(x^{*}\). A column shorter than \(\gamma\) cannot carry a nonzero coordinate of any minimizer.
\end{summary}
